% Source: https://artofproblemsolving.com/community/c3h2959098_geometry_circles
% Stable identity: aops-topic-2959098-q01
Triangles $PAC$ and $PDA$ share $\angle P$, and $\angle PAC=\angle PDA$ (tangent-chord angle), so they are similar: $\dfrac{AC}{DA}=\dfrac{PA}{PD}$. Likewise $\dfrac{BC}{DB}=\dfrac{PB}{PD}$.
\begin{asy}
usepackage("amsmath");
size(4.5cm);
defaultpen(fontsize(10pt));
pair O=(0,0);real R=1;pair P=(2.6,0);real t=degrees(acos(R/2.6));pair A=dir(t),B=dir(-t);
pair D=dir(165);path sec=P--D;pair C=intersectionpoints(P--D,circle(O,R))[0];
draw(circle(O,R));draw(P--A);draw(P--B);draw(P--D);draw(A--C--B);draw(A--D--B);
dot(P);label("$P$",P,E);label("$A$",A,N);label("$B$",B,S);label("$C$",C,NE);label("$D$",D,W);
\end{asy}

Since $PA=PB$, $\dfrac{AC}{DA}=\dfrac{BC}{DB}$, so $BC=\dfrac{AC\cdot DB}{DA}=\dfrac{4.5\cdot5}{6}=\boxed{3.75}$ cm.
